\scanpage{19}
\mspage{9}

Dualement, \deleted{\unsure{$A$ esp de|$A$ app de}} la donnée d'un carré-arêtes $A$ s'interprète comme la donnée d'un torseur sous un groupe $\mathbb F_2^D$, ce qui équivaut à la donnée d'une famille de \inserted{$\mathbb F_2$-}torseurs $(S_d)_{d\in D}$, \deleted{\unsure{indexés sur $S$}} \inserted{par}
\[
A\simeq\prod_{d\in D}S_d\qquad(\simeq\text{sections de }S\text{ sur }D)
\]
-- la donnée de cette famille équivalent à la donnée d'un ensemble (à quatre éléments) $S$ et d'un morphisme de degré $2$
\[
S\mathrel{\unsure{\text{---}|\longrightarrow}}D
\]
qui avec la formule précédente donne
\[
A\hookrightarrow\mathscr P_2(S)
\]
et la relation d'incidence entre $S$ et $A$\ldots

\uline{Formulaire} On a, pour un carré $Q$, les trois ensembles à deux éléments
\[
\begin{array}{ccc}
 C,&D,&\underline\omega\\
 \text{codiagonales}&\text{diagonales}&\text{orientations}
\end{array}
\]
plus les deux paquets de deux ens. à deux éléments (donc $4$ autres ens. à deux éléments)
\[
\begin{aligned}
S\xrightarrow{\mathrm{deg.}\,2}D,\qquad & A\xrightarrow{\mathrm{deg.}\,2}C\\
(D\simeq S/\underline a_S).\qquad &(C\simeq A/\underline a_A)
\end{aligned}
\]
On a le formulaire d'iso can.

% archive sheet number: 18

